% Original: PDF strana 9, donji slajd.
\slajdid{02-s018}
\begin{frame}[label=02-s018]{Indeksi potpunih proizvoda i zbirova}
\setlength{\abovedisplayskip}{4pt}\setlength{\belowdisplayskip}{4pt}
% element: 02-s018:e01
Oba oblika funkcije možemo jednostavnije zapisati indeksima proizvoda i zbirova.\par\vspace{1mm}
\begin{columns}[c,onlytextwidth]
\begin{column}{.58\textwidth}
% element: 02-s018:e02
Stanja ulaznih promenljivih tretiramo kao binarne cifre reda tabele. \alert{Indeks je decimalna vrednost tog binarnog broja.}
\par\vspace{1mm}
\[\begin{aligned}
\textcolor{DEcrvena}{F}&\textcolor{DEcrvena}{=\sum\nolimits_{C,B,A}(3,5,6)}\\[2mm]
F&=\prod\nolimits_{C,B,A}(0,1,2,4,7)
\end{aligned}\]
\par\vspace{1mm}
Cifra \(C\) je najveće težine, zatim cifra \(B\), pa cifra \(A\).
\end{column}
\begin{column}{.36\textwidth}\centering
\Oznaka{Indeksi i funkcionalna tabela}\vspace{2mm}
{\fontsize{9}{10}\selectfont\begingroup
\setlength{\tabcolsep}{3pt}
\setlength{\aboverulesep}{0pt}\setlength{\belowrulesep}{0pt}
\renewcommand{\arraystretch}{1.0}
\par\noindent
\begin{tabularx}{\linewidth}{l*{4}{>{\centering\arraybackslash}X}}
\toprule indeks&C&B&A&F\\\midrule
0&0&0&0&0\\\cmidrule[.2pt](lr){1-5}
1&0&0&1&0\\\cmidrule[.2pt](lr){1-5}
2&0&1&0&0\\\cmidrule[.2pt](lr){1-5}
3&0&1&1&1\\\cmidrule[.2pt](lr){1-5}
4&1&0&0&0\\\cmidrule[.2pt](lr){1-5}
5&1&0&1&1\\\cmidrule[.2pt](lr){1-5}
6&1&1&0&1\\\cmidrule[.2pt](lr){1-5}
7&1&1&1&0\\\bottomrule
\end{tabularx}\par
\endgroup
}
\end{column}
\end{columns}
\par\vspace{1mm}
% element: 02-s018:e03
\Oznaka{Vrednost binarnog broja \(b_{n-1}b_{n-2}\ldots b_0\)}
\[\textcolor{DEcrvena}{B=b_{n-1}2^{n-1}+b_{n-2}2^{n-2}+\cdots+b_1 2^1+b_0 2^0}\]
\end{frame}
